\chapter{Files Supplied}
\label{app:files}

\section{The program}

All seventeen must be in one directory on one device.

\begin{cctable}{File}{What it is}{1.7in}{3.12in}
\fname{CMDRCALC.PRG} & The program. This is the one you load. \\
\fname{OVL1.BIN} & The rest of the program. \ccprod{} is larger than the
machine's memory, so it keeps most of itself on the card and reads in whichever
part it needs at the time. Which file holds what is the program's business; all
you need to know is that every one of them has to be there. \\
\fname{OVL2.BIN} & \\
\fname{\ \ \ }\emph{\ldots{} through \ldots} & \\
\fname{OVL15.BIN} & \\
\fname{OVL16.BIN} & \\
\end{cctable}

About 135K in total. The byte counts are from one build and move as the
program is worked on; what matters is that all seventeen files are present.

\section{Example workbooks}

\subsection{In Commander Calc's own format}

Open these with \key{F3}.

\cccmd[]{INVOICE.X16S}

One worksheet. A small invoice: currency, a date, a computed line total from
\fml{=B5*C5}, a \fn{SUM} of the lines, VAT as a percentage, and a total due.

Shows what the currency, date and percentage formats look like --- and, since
none of them can be set from the keyboard, it is the shortest way to see them
at all.

\cccmd[]{RAIN.X16S}

Two worksheets, \emph{Readings} and \emph{Summary}. Twelve months of rainfall
on the first; the second reads across the sheet boundary with
\fml{=SUM(Readings!B2:B13)}, \fn{MAX}, \fn{MIN}, \fn{AVERAGE}, \fn{COUNT} and
an \fn{IF}.

The one to look at if cross-sheet references are what you are trying to
understand.

\cccmd[]{KINDS.X16S}

One worksheet, one row per kind of thing a cell can be: text, a number, a
negative number, a very small number, currency, a percentage, a fixed three
places, a date, a boolean, \fml{=2\^{}10} which should read 1024, and
\fml{=1/0} which should read \errv{\#DIV/0!} and does.

A reference card for what everything looks like on screen.

\cccmd[]{NDQNAT.X16S}

Two worksheets: five hundred days of real NASDAQ Composite prices, and a
summary. Large.

This one is a stopwatch. It is exactly the same data as the \fname{.XLSX}
version, and the point is the difference in how long the two take to open ---
seventeen seconds against fifty-one.

\subsection{Excel workbooks}

Open these with \key{F7}.

\cccmd[]{BUDGET.XLSX}

One worksheet. Money, percentages and totals across a ten-row table of
categories, using currency and percentage formats, \fn{SUM}, \fn{MAX},
\fn{AVERAGE}, subtraction and division.

\cccmd[]{STOCK.XLSX}

One worksheet, deliberately in no order at all. Mixed text and number columns,
line values from \fml{=C*D}, \fn{SUM} and \fn{COUNT}.

Made for trying \menupath{Data > Sort up} on: put the cursor in any column, on
the first row of data, and watch whole rows move. It is also the best of these
for trying the \menupath{Chart} menu against.

\cccmd[]{TOTALS.XLSX}

Three worksheets: \emph{Q1}, \emph{Q2} and \emph{Summary}. The third reads the
first two --- \fml{=Q1!B5}, \fml{=MAX(Q1!B2:B4,Q2!B2:B4)} --- so it also
demonstrates that recalculation crosses worksheets.

\cccmd[]{FUNCS.XLSX}

One worksheet with one row per function and the expected answer written beside
it. Every one of the fourteen, plus the \lit{\^{}} operator, a comparison, and
\fml{=A20/0} producing \errv{\#DIV/0!} on purpose.

The quickest way to check that a build is working, and a usable crib sheet.

\begin{ccnote}
Every formula in these files is one \ccprod{} can compile. They are not test
fixtures --- they are small spreadsheets built to be used, each exercising a
different part of the program.
\end{ccnote}

\section{Files the program creates}

\begin{cctable}{File}{When}{1.6in}{3.22in}
\fname{CMDRCALC.TMP} & During every save. Deleted, or renamed over your
workbook, when the save finishes. If you find one lying about, a save was
interrupted --- your previous workbook is intact and the temporary file can be
deleted. \\
\fname{CHART.BMP} & When \key{S} is pressed on a chart. Always that name, so
each one overwrites the last. Chapter~\ref{ch:charting}. \\
\end{cctable}
